\section{関連研究}
\label{sec_related}
画像とユーモアを扱う研究では，画像に付けた一言の面白さを理解する課題\cite{ref_newyorker}や，ボケてを含む大喜利のデータから回答を生成する課題\cite{ref_oogirigo}が提案されている．
これらの研究は，画像と回答の組を一様に扱うことが多い．
しかし，お題画像には画像中の題に答えるものと画像そのものに一言を付けるものが混在しており，両者を区別しなければ，モデルが画像中の文字を読む必要があるかを評価できない．

画像中の文字を読む課題としては，情景の中の文字を手がかりに質問に答える Visual Question Answering の課題\cite{ref_textvqa,ref_stvqa}が知られている．
これらの課題では，看板や商品の表示のように，写っている世界の中の文字を読むことが主な対象である．
一方，お題画像には，写真の上に重ねた字幕や，画像の外の帯に書いた文が多く含まれ，文字がどこに書かれているかが，その文字の役割を推し量る手がかりになる．
本稿の層の属性は，この違いを明示的に記録するものである．

ラベルの付与では，判断者の間で正解が揺れることが知られており，揺れを誤りとして除くのではなく，データの性質として扱うべきとする議論がある\cite{ref_hlv}．
本稿で扱うお題文の判断の揺れは，判断者の解釈の違いではなく，画像から観測できない投稿者の意図を定義が要求することに由来する．
本稿では，観測できる属性だけでラベルを定義し，解釈に依存する判断を属性からの規則に移すことで，この揺れを扱う．

VLM の判定を高速な分類器に移す方法は，知識蒸留\cite{ref_kd}として知られる．
本研究のパッケージも，Qwen3-VL\cite{ref_qwen3vl}を教師とし，SigLIP 2\cite{ref_siglip2}の埋め込みに基づく生徒を用いる．
本稿では，生徒の入力に日本語の文字認識 (Optical Character Recognition，OCR) の YomiToku\cite{ref_yomitoku}による書き起こしを加えた場合も比べる．
